\section{Introduction}

Vision-Language-Action (VLA) models have demonstrated strong capabilities
in robotic manipulation. However, their computational cost and inference
latency can limit their use in real-time robotic systems. Improving
inference efficiency while preserving action quality is therefore important
for deploying VLA models on practical robot platforms.

Existing approaches improve VLA efficiency at either the model or system
level. Model-level methods include fast--slow architectures, lightweight
vision-language backbones, KV caching, and early-exit strategies. System-level
methods include quantization, CUDA Graphs, and custom CUDA operators.
These approaches reduce different sources of latency, but action generation
based on iterative denoising can remain computationally expensive.

In this technical report, we investigate a one-step action-generation head
based on [define ``drifting'' here]. The proposed head replaces the
multi-step denoising procedure of the original diffusion-transformer
(DiT) action head with a single generation step. In our experiments on
NVIDIA GR00T N1.7, this modification reduces [action-head/end-to-end]
inference latency to approximately 10\% of the original value under the
same evaluation conditions. The head is designed as a modular modification
to the VLA architecture, although its applicability to other models remains
to be evaluated.

One-step generation also creates a consistency problem during asynchronous
action prefetching: a newly predicted action chunk may not account for
actions that have already been committed to the robot's control queue.
To address this issue, we introduce an overlap-aware action head inspired
by training-time RTC [define and cite RTC]. It conditions generation on
the committed action prefix and preserves valid prefix coordinates exactly.

The main contributions of this report are:
\begin{itemize}
    \item We implement and evaluate a one-step action-generation head for
    NVIDIA GR00T N1.7.
    \item We describe the consistency limitation of combining one-step
    generation with asynchronous action prefetching.
    \item We introduce an overlap-aware formulation that represents the
    committed action prefix, coordinate-wise validity, overlap length, and
    estimated execution offset while preserving valid prefix coordinates.
\end{itemize}
